\chapter{Completing the run argument with an explicit logarithm bound}
{\small\sffamily Source: \nolinkurl{repository/collatz_reconstruction/research_program/defect_rank_descent_20260915/logarithmic_run_exclusion.md}.\par}


15 September 2026. This is part of the
\nolinkurl{defect_rank_descent_20260915} continuation. The original map,
word order, integral sections and first-jet differentials are those of
\nolinkurl{note.md}. This companion proves two additional all-length
results. It does not assert an all-source Collatz conclusion.

\section{1. Established transcendence input and the exact
specialization}

We use the rational-number specialization of Matveev's lower bound, as
printed in Theorem 2.1 of Languasco, Luca, Moree and Togbe,
\emph{Sequences of integers generated by two fixed primes}, Abhandlungen
aus dem Mathematischen Seminar der Universitat Hamburg 95 (2025),
123--148, DOI \nolinkurl{10.1007/s12188-025-00293-9}. The complete
theorem statement and its logarithmic-height convention were read in the
publisher's HTML. That paper cites Matveev's original theorem and the
Bugeaud--Mignotte--Siksek formulation. This external transcendence
theorem is an explicit proof dependency, not a theorem claimed to have
been proved by our finite checker.

In that statement, for positive rational numbers gamma\_i, nonzero
integers b\_i, and nonzero \nolinkurl{Lambda=product gamma_i^b_i-1},

\SourceMath{\log|\Lambda|>
-1.4\,30^{k+3}k^{4.5}(1+\log B)\prod_i h(\gamma_i),
\quad B\ge\max_i|b_i|.}{L1}

Here \nolinkurl{h(r/s)=max(log|r|,log s)} for reduced positive
rationals. We use only \nolinkurl{gamma_1=2,gamma_2=3}, with exponents
\nolinkurl{A,-m}, so their heights are log 2 and log 3. The original
coefficient expression is

\SourceMath{\Lambda=2^A3^{-m}-1>0,
\quad A>m\ge1,
\quad B=A.}{}

Unique prime factorization proves it nonzero. No exponent in this
application is zero. The coefficient in (L1) has the rigorous rational
upper bound

\SourceMath{1.4\,30^5\,2^{4.5}\log2\log3
<\frac{14}{10}30^5\,24\frac7{10}\frac{11}{10}
=628689600<10^9.}{}

We used \nolinkurl{sqrt(2)<3/2}, \nolinkurl{1/2<log2<7/10}, and
\nolinkurl{1<log3<11/10}. The latter inequalities follow from the
explicit series bounds in Section 2. Put \nolinkurl{C_*=10^9}. Thus the
only transcendence estimate used below is

\SourceMath{\boxed{\log\Lambda>-C_*(1+\log A).}}{L2}

We retain the original A,m in every substitution. This does not identify
a coefficient contraction with an actual descent. The additive
correction and integer source congruences enter separately through the
run argument in \nolinkurl{note.md}.

\section{2. A fully rational interval excluding an infinite collection
of approximations}

For an integer x\textgreater=2, set \nolinkurl{z=(x-1)/(x+1)}.
Integration of the geometric series, or its finite identity with
remainder, gives

\SourceMath{2\sum_{j=0}^{N-1}\frac{z^{2j+1}}{2j+1}
<\log x
<2\sum_{j=0}^{N-1}\frac{z^{2j+1}}{2j+1}
+\frac{2z^{2N+1}}{(2N+1)(1-z^2)}.}{L3}

All omitted terms are positive, and replacing their denominators by
\nolinkurl{2N+1} bounds their sum by the last geometric tail. The
implementation uses N=128 for x=2 and x=3. Write the resulting rational
intervals as \nolinkurl{(l_2,u_2)} and \nolinkurl{(l_3,u_3)}. Then

\SourceMath{\gamma_-:=l_3/u_2<\gamma:=\log3/\log2<\gamma_+:=u_3/l_2.}{}

The following certificate consists entirely of integer and rational
comparisons:

\SourceMath{\frac{83130157078217}{52449289519716}
<\gamma_-<\gamma_++2^{-128}
<\frac{350861368503572}{221368876767703}.}{L4}

Its two endpoint numerators and denominators satisfy

\SourceMath{350861368503572\cdot52449289519716
-83130157078217\cdot221368876767703=1,}{}
\SourceMath{52449289519716+221368876767703
=273818166287419>226\cdot10^{12}.}{L5}

The final endpoint pair is obtained by a mediant search whose steps and
hash are recorded; the proof needs only (L4)--(L5), not confidence in
that search.

\subsection{Elementary denominator lemma}

For reduced fractions l=p/q\textless h=r/s with rq-ps=1, every rational
x/y strictly between them, with x,y positive integers, satisfies
y\textgreater=q+s.

\textbf{Proof.} The integers \nolinkurl{k=xq-py} and \nolinkurl{j=ry-sx}
are both positive. Direct expansion gives \nolinkurl{y=s k+q j}, hence
y\textgreater=q+s. This holds even when x/y has not been reduced. QED.

Applying this lemma proves

\SourceMath{\boxed{
0<A/m-\gamma<2^{-128},\quad 1\le m\le226\cdot10^{12}
\quad\text{has no integral solution }(A,m).
}}{L6}

The map from the original period coordinates to the rational number is
\nolinkurl{(A,m)->A/m}. Its fibres consist of integer multiples of a
reduced numerator/denominator pair. The denominator lemma applies
directly to every unreduced member; no original multiplicity is lost.

\section{3. The contracting pure-two comparison is always strictly
separated}

Recall

\SourceMath{D_2(a,b)=2^{a+2b}-3^{a+b},\qquad a,b\ge1.}{}

\subsection{Theorem E. Complete positive-gap inequality}

On the complete integer domain \nolinkurl{D_2(a,b)>0},

\SourceMath{\boxed{D_2(a,b)>3^a-2^a.}}{L7}

\textbf{Proof.} The elementary run-cone proof in \nolinkurl{note.md}
establishes (L7) for all \nolinkurl{b>=ceil(3a/2)}. Consider an alleged
remaining pair with \nolinkurl{0<D_2<=3^a-2^a}. It has

\SourceMath{a<b<\lceil3a/2\rceil,\quad
m=a+b<5a/2,\quad A=a+2b<4a.}{L8}

The lower inequality b\textgreater a follows from
\nolinkurl{(2/3)^a(4/3)^b <= (8/9)^a<1} when b\textless=a. The integer
endpoint for b gives the strict A\textless4a inequality in both parity
cases.

The original relative gap obeys

\SourceMath{0<\Lambda=D_2/3^{a+b}<3^{-b}<3^{-a}.}{}

Together with (L2) and log 3\textgreater1, this yields

\SourceMath{a<C_*(1+\log(4a)).}{L9}

It forces a\textless10\^{}12. For an explicit check, the positive finite
sum \nolinkurl{sum_(j=0)^100 36^j/j!} exceeds \nolinkurl{452*10^12}, so
\nolinkurl{log(4*10^12)<36}. At a=10\^{}12, the left minus right of (L9)
is positive; its derivative \nolinkurl{1-C_*/a} is positive thereafter.
These inequalities are checked by rational arithmetic.

For a\textgreater=128, therefore m\textless226*10\^{}12 and

\SourceMath{0<A/m-\gamma
=\frac{\log(1+\Lambda)}{m\log2}
<2\,3^{-b}/m
<2\,3^{-128}<2^{-128}.}{}

This contradicts (L6). For the remaining 1\textless=a\textless=127, the
finite certificate computes the least b with D\_2\textgreater0 and
verifies (L7) there. Increasing b sends D\_2 to
\texttt{3D\_2+2\^{}\{a+2b\}}, so the inequality persists. These 127
exact rows cover all remaining b, rather than truncating b at a search
bound. QED.

\subsection{Corollary E1. Strict original descent at the exact
comparison threshold}

For every actual positive odd source whose word is
\nolinkurl{1^a,c_1,...,c_b}, with a,b\textgreater=1 and every
c\_i\textgreater=2, the original endpoint is strictly below the original
source throughout the domain

\SourceMath{\boxed{2^{a+2b}>3^{a+b}.}}{L10}

\textbf{Proof.} For a pure-two tail, the exact source parameter and
displacement from \nolinkurl{note.md} are

\SourceMath{3^a t-1=4^b z,\quad z\ge1,
\qquad n-y=2\left(\frac{D_2z+2^a}{3^a}-1\right).}{}

Equation (L7) makes the inner integer at least two, so
n-y\textgreater=2. A tail containing a higher exponent is handled by the
proved monotone affine comparison in Lemma C2 of \nolinkurl{note.md}, on
precisely the same original source. QED.

This is stronger than the elementary 3/2 cone and the certified
2124/1507 cone. Those earlier arguments remain useful as shorter
self-contained certificates that do not use (L1). The new condition
refers to the pure-two COMPARISON multiplier; it does not assert that
every arbitrary word with contracting actual multiplier descends. The
actual higher exponents, offsets, source congruences, and intermediate
integers remain attached. A zero displacement left possible by the
earlier weak pure-two lemma is now excluded on its actual positive
domain by (L7), not deleted from its source declaration.

\section{4. A global restriction on the rise beginning at an actual
cycle minimum}

Let an actual nontrivial positive cycle have primitive odd-return length
m, original exponent sum A, and least original vertex n.~It cannot have
exponent \textgreater=2 leaving n, because that would give
\nolinkurl{(3n+1)/4<n}. Its first a\textgreater=1 exponents are
therefore a maximal run of ones. The exact original source formula gives

\SourceMath{n=2^{a+1}t-1,\qquad t>0\text{ odd}.}{L11}

The common source cover in \nolinkurl{note.md} gives
n\textgreater=s\_0=1166401. In the complete original source of this run,
the least allowed vertex is more precise than the maximum of two
inequalities. It is

\SourceMath{t_a=\text{least positive odd integer at least }(s_0+1)/2^{a+1},
\qquad s_a=2^{a+1}t_a-1.}{L12}

This is an exact progression intersection, with inverse
\texttt{t=(n+1)/2\^{}\{a+1\}}. For example at a=19, (L12) gives 3145727,
not the weaker bound 1166401. At a=17 it gives 1310719. The certificate
retains every t\_a and s\_a.

There is further source information in the already certified
supported-zero fibre. The complete equation support E\_B in
\nolinkurl{note.md} has its actual finite integral contraction to 1. A
nontrivial cycle cannot contain a vertex of E\_B. Define
\nolinkurl{tilde t_a} to be the least odd t\textgreater=t\_a for which
\nolinkurl{2^(a+1)*t-1} is not in E\_B, and put

\SourceMath{\widetilde s_a=2^{a+1}\widetilde t_a-1.}{L12a}

This exists because E\_B is finite. For a\textless=127 the delivered
\nolinkurl{retained_zero_minima.json} lists every skipped original
parameter and its finite integral first-jet witness. There are only nine
skipped sources, with 639 original returns in their witness expansions.
Every one is independently checked to belong to the previously
constructed E\_B with its recorded height. Thus this reuses the retained
source behind a supported zero; it is not a new sweep of starting
integers or an assertion that a zero value determines its preimage
without the map.

For example, at a=17 the original value 1310719 is in E\_B and has
height 105, although it lies above the contiguous interval bound. Its
exclusion makes the next allowed parameter 7 and the next minimum bound
1835007. The unweighted equations and the finite first-jet witness
behind that exclusion remain in the record. For larger a the inequality
\nolinkurl{widetilde s_a>=2^(a+1)-1} suffices; no unknown infinite orbit
is queried.

\subsection{Theorem F. All-length minimum-run exclusion}

Every actual nontrivial positive cycle satisfies

\SourceMath{\boxed{m>226a,}}{L13}

where a is the maximal initial exponent-one run at its actual minimum.
The coefficient 226 is a proved workbench bound, not a claimed optimal
constant or published record.

\textbf{Proof.} Assume instead m\textless=226a. Multiplying the original
cycle equations gives

\SourceMath{0<\delta:=A\log2-m\log3
=\sum_{i=1}^m\log\left(1+\frac1{3x_i}\right)
\le\frac{m}{3\widetilde s_a}.}{L14}

All x\_i\textgreater=3, so \nolinkurl{2^A<= (10/3)^m<4^m}, hence
A\textless2m\textless=452a. Also Lambda=e\^{}delta-1 is precisely the
positive expression in (L2).

First let a\textgreater=128. Equation (L11) gives
\texttt{delta\textless{}=226a/{[}3(2\^{}\{a+1\}-1){]}\textless{}1/2}.
The last bound is checked at 128 and persists:
\texttt{a/(2\^{}\{a+1\}-1)} decreases, as direct cross-multiplication
shows. For \nolinkurl{0<delta<1/2}, integration of e\^{}x and
\texttt{e\^{}\{1/2\}\textless{}2} give \nolinkurl{e^delta-1<2delta}.
Consequently

\SourceMath{\Lambda<226a\,2^{-a}.}{}

Combining with (L2) gives

\SourceMath{a\log2<C_*(1+\log(452a))+\log(226a).}{L15}

Since log2\textgreater1/2, the rational certificate at a\_0=10\^{}12,

\SourceMath{\log(452a_0)<36,\qquad a_0/2>37C_*+36,
\quad a_0>2(C_*+1),}{}

and the derivative \nolinkurl{1/2-(C_*+1)/a} prove a\textless a\_0. Thus
m\textless226*10\^{}12. On the other hand (L14) gives

\SourceMath{0<A/m-\gamma\le\frac1{3\widetilde s_a\log2}<2^{-a}\le2^{-128},}{}

contradicting (L6).

It remains to cover 1\textless=a\textless=127. For each such a, the
certificate constructs a rational bracket
\nolinkurl{p_a/q_a < gamma_- < gamma_+ < r_a/s'_a} with determinant one,
denominator sum greater than 226a, and

\SourceMath{\gamma_++\frac1{3\widetilde s_a l_2}<\frac{r_a}{s'_a}.}{L16}

The exact table contains all 127 source parameters, bounds and endpoint
pairs. Equation (L14) places A/m strictly within this bracket. The
denominator lemma gives m\textgreater226a, the desired contradiction.
All rows are checked with exact fractions and the same certified
logarithm intervals. QED.

Retaining the original odd parameter first improves the weaker
interval-only ratio 191 to 213. The nine original source exclusions in
(L12a), already supported by E\_B, then improve 213 to 226. The checker
verifies both the plain-progression baseline and the support-refined
table. No additional discovery source was iterated for this improvement.
The previously retained zero fibre is used with its actual original
source and boundary witness, not as a bare zero symbol.

\subsection{Corollary F1. A complete ordering family of cycles is
excluded}

There is no nontrivial positive cycle whose primitive original exponent
word has all its exponent-one letters in a single cyclic run. The other
exponents may be arbitrarily large and their number is unbounded.

\textbf{Proof.} Such a cycle has, after its actual minimum, a
consecutive ones and then b\textgreater=1 exponents \textgreater=2. The
ones strictly increase each vertex, while every later edge strictly
decreases its vertex above 1, so the beginning of the run is indeed the
actual minimum and m=a+b. The exponent sum satisfies
A\textgreater=a+2b=2m-a. The same original cycle product gives

\SourceMath{(6/5)^m\le2^a.}{}

Since \nolinkurl{(6/5)^4>2}, this forces m\textless4a, contradicting
(L13). A word with no ones decreases at every nonbase vertex and cannot
be such a cycle; a word with only ones increases at every vertex. Thus
no excluded edge case is assigned a positive periodic source. QED.

This result excludes an all-length family of actual period-torsion
modules through their original arithmetic realizations. It does not
assert that arbitrary abstract circles with this word pattern are
absent. Their rational affine primitives, integral forcing obstructions
and supported zero values remain available; what has been excluded is a
positive integral cycle realizing that pattern.

\section{5. Exact relationship with the first-jet object}

No new cohomological quotient is introduced. Every path certified by
Corollary E1 uses the same original equations and satisfies

\SourceMath{d_\epsilon(\epsilon B_p(n))=\epsilon V_n-\epsilon V_{g_p(n)}.}{}

Since the endpoint is lower, an existing integral endpoint witness
splices to the source witness. Original cycle modules have first-jet
index equal to their ORIGINAL primitive period. Theorem F and Corollary
F1 restrict the original words and original minimum placements that can
produce such a module. They do not infer disappearance merely from a
zero scalar or from passing to rational coefficients.

The proven source-cover bound s\_0 in Theorem F is a finite arithmetic
dependency verified in this same package. The transcendence estimate
(L1) is a sourced established theorem. Equations (L3)--(L16) and the
full run proofs specify every intervening inequality. No probability
law, unproved stopping-time estimate, or assumed common nilpotent
exponent is used.

\section{6. Execution and remaining source}

\nolinkurl{logarithm_certificate.py} uses only exact rational
arithmetic. It recomputes the positive series bounds, the large Farey
bracket and its determinant, the 127 small pure-run rows, all 127 exact
minimum-run source brackets, and the nine reused original supported-zero
witnesses. \nolinkurl{logarithm_verification.json} records the complete
output. \nolinkurl{verify.py} includes this certificate and rejects
changed bracket endpoints and an invalid coefficient claim. No numerical
logarithm or unchecked floating-point continued fraction is a proof
premise.

The finite rational certificate does not independently verify Matveev's
theorem. The theorem is explicitly sourced; the subsequent finite steps
and arithmetic instantiation are checkable here. These all-length
results remain additional to, not a rebranding of, the finite-source
cover and its 1506-exponent-one-count exclusion.

The actual remaining cycle source has at least two cyclic exponent-one
runs and obeys the original minimum-run restriction. The aperiodic
source still requires control of genuine packet switches and their
resonant source fibres. The nonzero-defect rank in \nolinkurl{note.md}
handles repeated packets, not arbitrary switching. Neither new argument
proves all such sources converge.

\section{Sources}

{[}M{]} E. M. Matveev, \emph{An explicit lower bound for a homogeneous
rational linear form in the logarithms of algebraic numbers. II},
Izvestiya: Mathematics 64 (2000), 1217--1269. DOI
\nolinkurl{10.1070/IM2000v064n06ABEH000314}. The original article's
bibliographic record was checked; its full proof was not re-audited in
this session.

{[}LLMT{]} A. Languasco, F. Luca, P. Moree and A. Togbe, \emph{Sequences
of integers generated by two fixed primes}, Abhandlungen aus dem
Mathematischen Seminar der Universitat Hamburg 95 (2025), 123--148,
Theorem 2.1 and the immediately preceding height definition. DOI
\nolinkurl{10.1007/s12188-025-00293-9}. Full readable source:
\nolinkurl{https://link.springer.com/article/10.1007/s12188-025-00293-9}.
This is the inspected formulation of {[}M{]} used in (L1).
